% =========================================================================
% BUKIT — Memoria del TFG (estructura según la plantilla de la Escuela de
% Ingeniería de Bilbao + guía de TFG de la UPV/EHU)
% Documento raíz: aquí SOLO se ensambla el documento. El contenido vive
% en los módulos (config/, frontmatter/, chapters/, appendices/).
%
% Compilar:  desde VSCode al guardar (LaTeX Workshop), o `latexmk` en esta
% carpeta (requiere Perl). El PDF queda en build/main.pdf.
% =========================================================================

% Clase provisional: el formato oficial (Calibri, márgenes, títulos...)
% se aplicará en config/style.tex en la fase de diseño.
\documentclass[12pt, a4paper, oneside]{report}

% --- Configuración global (el orden importa) ---
% =========================================================================
% Variables globales del documento.
% Cambiar un dato AQUÍ lo actualiza en todo el documento (portada,
% metadatos del PDF...). No escribir nunca estos datos "a mano" en otros
% archivos: usar siempre estos comandos.
% =========================================================================

% --- Trabajo ---
\newcommand{\thesisTitle}{BUKIT: red social orientada a objetivos}% TODO: título definitivo
\newcommand{\thesisSubtitle}{Memoria del Trabajo de Fin de Grado}%
\newcommand{\thesisDegree}{Grado en Ingeniería Informática de Gestión y Sistemas de Información}%
\newcommand{\thesisCourse}{2026--2027}% curso académico
\newcommand{\thesisDate}{Bilbao, \today}% TODO: <Lugar, día, mes, año> de la entrega/defensa
\newcommand{\thesisYear}{2027}% año de elaboración y defensa

% --- Personas (la plantilla del centro usa "Apellido1 Apellido2, Nombre") ---
\newcommand{\thesisAuthor}{Liviu Deleanu}% nombre natural (para metadatos del PDF)
\newcommand{\thesisAuthorFormal}{Deleanu, Liviu}% TODO: confirmar formato con dos apellidos
\newcommand{\thesisDirector}{Villamañe Gironés, Mikel}%
\newcommand{\thesisCoDirector}{Apellido1 Apellido2, Nombre}% TODO: solo si hay codirector/a

% --- Universidad ---
\newcommand{\thesisUniversity}{Universidad del País Vasco / Euskal Herriko Unibertsitatea}%
\newcommand{\thesisFaculty}{Escuela de Ingeniería de Bilbao}%

% --- Palabras clave del resumen trilingüe ---
\newcommand{\thesisKeywordsEs}{red social, objetivos, monorepo, TypeScript, PostgreSQL}% TODO
\newcommand{\thesisKeywordsEu}{sare soziala, helburuak, monorepo, TypeScript, PostgreSQL}% TODO: revisar euskera
\newcommand{\thesisKeywordsEn}{social network, goals, monorepo, TypeScript, PostgreSQL}% TODO
   % Variables globales (título, autor, fechas...)
% =========================================================================
% Carga y configuración de paquetes, agrupados por función y con el porqué
% de cada uno. hyperref se carga casi al final (requisito del paquete).
% =========================================================================

% --- Codificación y fuentes (pdflatex) ---
\usepackage[utf8]{inputenc}   % Codificación de entrada UTF-8
\usepackage[T1]{fontenc}      % Acentos correctos en el PDF (copiables/buscables)

% --- Idiomas: español principal; euskera e inglés para el resumen trilingüe ---
% es-tabla: "Tabla" en vez de "Cuadro".
% es-noshorthands: desactiva los caracteres activos del español (no los
%   necesitamos: escribimos los acentos en UTF-8) y evita un conflicto real
%   entre babel-spanish y el hook de biblatex-apa sobre \@makecaption
%   ("Generic hooks cannot be added to '\@makecaption'").
\usepackage[english, basque, spanish, es-tabla, es-noshorthands]{babel}
\usepackage{csquotes}         % Comillas correctas según idioma (lo pide biblatex)
% Nomenclatura de la guía de la EIB: "Anexo" para los apéndices e "Índice de
% ilustraciones" para las figuras ("Índice de tablas" ya es el defecto con
% es-tabla):
\addto\captionsspanish{%
  \renewcommand{\appendixname}{Anexo}%
  \renewcommand{\listfigurename}{Índice de ilustraciones}%
}

% --- Gráficos y tablas ---
\usepackage{graphicx}
\graphicspath{{figures/}}     % Las imágenes se buscan siempre en figures/
\usepackage{booktabs}         % Tablas profesionales (\toprule, \midrule, \bottomrule)

% --- Utilidades ---
\usepackage{xspace}           % Espaciado correcto tras macros de texto (\bukit)
\usepackage{xcolor}           % Colores (usado por \todonote)
\usepackage{acronym}          % Lista de acrónimos (ver frontmatter/acronyms.tex)

% --- Bibliografía: biblatex + biber, en normas APA (exigido por la guía) ---
% Citar con \parencite{clave} -> "(Williams, 1995)"
%       o  \textcite{clave}   -> "Williams (1995)"
\usepackage[backend=biber, style=apa]{biblatex}
\DeclareLanguageMapping{spanish}{spanish-apa}
% biblatex-apa no trae localización en euskera (basque-apa.lbx); mapeamos a
% la española para silenciar el aviso. Solo afectaría a citas hechas DENTRO
% del bloque en euskera (el Laburpena), donde nunca citamos.
\DeclareLanguageMapping{basque}{spanish-apa}

% --- Enlaces y metadatos del PDF (cargar al final) ---
\usepackage[hidelinks]{hyperref}
\hypersetup{
  pdftitle={\thesisTitle},
  pdfauthor={\thesisAuthor},
  pdfsubject={\thesisSubtitle},
  pdfkeywords={\thesisKeywordsEs},
}
   % Paquetes y su configuración
% =========================================================================
% Comandos propios. Definir aquí cualquier cosa que se repita en el texto:
% si mañana cambia cómo se escribe, se cambia en un único sitio.
% =========================================================================

% Nombre del producto, siempre con el mismo formato.
\newcommand{\bukit}{\textsc{Bukit}\xspace}

% Términos en inglés, en cursiva de forma consistente: \eng{backend}
\newcommand{\eng}[1]{\textit{#1}}

% Notas de trabajo, visibles en rojo en el PDF. Buscar "todonote" y
% eliminarlas todas antes de entregar.
\newcommand{\todonote}[1]{\textcolor{red}{\textbf{[TODO: #1]}}}
     % Comandos propios
% =========================================================================
% Formato y diseño del documento. PENDIENTE: se aplicará en la fase de
% diseño. Los requisitos ya conocidos (guía de TFG de la UPV/EHU) quedan
% anotados aquí para no perderlos:
%
%   - Márgenes: 2,5 cm por los cuatro lados        -> paquete geometry
%   - Interlineado sencillo; 6 pt entre párrafos   -> parskip / setspace
%   - Tipo de letra: Calibri                       -> con pdflatex no existe
%     Calibri; opciones: paquete `carlito` (clon métricamente compatible)
%     o compilar con XeLaTeX/LuaLaTeX + fontspec usando la Calibri de
%     Windows. Decidir en la fase de diseño.
%   - Títulos: capítulo MAYÚSCULA 14 pt negrita; apartado 1er nivel
%     MAYÚSCULA 12 pt negrita; 2º nivel minúscula 12 pt negrita; 3er nivel
%     minúscula 12 pt cursiva                      -> paquete titlesec
%   - Numeración de apartados: 1., 1.1., 1.1.1.    [ya es el defecto]
%   - Tablas/figuras: doble numeración capítulo.orden  [report ya lo hace]
%   - Impresión a doble cara                       -> \documentclass[..., twoside]
%   - Extensión objetivo: 40 páginas (máx. 60 excepcional)
%   - Identidad visual UPV/EHU en la portada (logos)
%
% Ejemplos de lo que vivirá aquí (dejar comentado de momento):
% \usepackage[a4paper, margin=2.5cm]{geometry}
% \usepackage{parskip}
% \usepackage{titlesec}
      % Formato y diseño (pendiente de definir)

% --- Bibliografía ---
\addbibresource{bibliography/references.bib}

\begin{document}

% ----- Portada -----
% =========================================================================
% Portada según la plantilla de la Escuela de Ingeniería de Bilbao:
% logos, "GRADO EN... / TRABAJO FIN DE GRADO", título, foto opcional,
% iconos de sostenibilidad, estudiante, director/a, curso y fecha.
% La estructura ya es la correcta y el formato visual global está aplicado
% en config/style.tex; en la portada solo quedan pendientes los logos
% UPV/EHU, la foto opcional y los iconos de sostenibilidad.
% =========================================================================

% La portada no genera ancla de página en el PDF: el entorno titlepage
% resetea el contador y, sin esto, hyperref avisa de dos páginas "1"/"i".
\hypersetup{pageanchor=false}
\begin{titlepage}

  % --- Logos institucionales (normativa de identidad visual UPV/EHU) ---
  % TODO: guardar los logos en figures/logos/ y descomentar:
  % \noindent
  % \includegraphics[height=1.6cm]{logos/upv-ehu}\hfill
  % \includegraphics[height=1.6cm]{logos/centro}
  \vspace*{1cm}

  {\noindent\large\bfseries \MakeUppercase{\thesisDegree}\par}
  \smallskip
  {\noindent\Huge\bfseries TRABAJO FIN DE GRADO\par}
  \vspace{1.5cm}

  \begin{center}
    {\LARGE\bfseries\itshape \thesisTitle\par}

    \vfill
    % TODO: foto representativa del trabajo (tamaño máximo: 10x10 cm):
    % \includegraphics[width=8cm]{portada-foto}
    \vfill
    % TODO: iconos de sostenibilidad (tamaño: 2x2 cm cada uno):
    % \includegraphics[height=2cm]{logos/ods-09}
  \end{center}
  \vfill

  \noindent\textbf{Estudiante:} \thesisAuthorFormal\par
  \medskip
  \noindent\textbf{Director/Directora:} \thesisDirector\par
  % TODO: descomentar solo si hay codirector/a:
  % \noindent\textbf{Codirector/Codirectora:} \thesisCoDirector\par
  \bigskip
  \noindent\textbf{Curso:} \thesisCourse\hfill\textbf{Fecha:} \thesisDate\par

\end{titlepage}
\hypersetup{pageanchor=true}


% ----- Preliminares (numeración romana), en el orden de la guía de la EIB:
%       resumen -> índices de tablas/ilustraciones -> acrónimos -> empleo de
%       IA -> índice de la memoria -----
\pagenumbering{roman}
% % =========================================================================
% Agradecimientos (opcional). Para incluirlos, descomentar su \input en
% main.tex; si no, este archivo simplemente no se usa.
% =========================================================================

\chapter*{Agradecimientos}

\todonote{Escribir agradecimientos, o dejar este archivo sin incluir.}
   % Opcional: descomentar si se incluyen
% =========================================================================
% Resumen trilingüe (castellano / euskera / inglés) + palabras clave,
% como exige la plantilla del centro. La guía además pide que recoja de
% forma sintética objetivos, líneas de trabajo y resultados (máx. 3 págs).
% Consejo: es lo ÚLTIMO que se escribe — resume el trabajo ya terminado.
% =========================================================================

\chapter*{Resumen}
\addcontentsline{toc}{chapter}{Resumen}

\todonote{Escribir el resumen al final del proyecto: objetivos, líneas
principales de trabajo y resultados más significativos.}

\bigskip
\noindent\textbf{Palabras clave:} \thesisKeywordsEs

% --- Euskera ---
\begin{otherlanguage}{basque}
  \chapter*{Laburpena}
  \addcontentsline{toc}{chapter}{Laburpena}

  \todonote{Laburpena proiektuaren amaieran idatzi.}

  \bigskip
  \noindent\textbf{Gako-hitzak:} \thesisKeywordsEu
\end{otherlanguage}

% --- Inglés ---
\begin{otherlanguage}{english}
  \chapter*{Abstract}
  \addcontentsline{toc}{chapter}{Abstract}

  \todonote{Write the abstract at the end of the project.}

  \bigskip
  \noindent\textbf{Keywords:} \thesisKeywordsEn
\end{otherlanguage}
             % Resumen trilingüe + palabras clave

\clearpage\phantomsection\addcontentsline{toc}{chapter}{\listtablename}
\listoftables
\clearpage\phantomsection\addcontentsline{toc}{chapter}{\listfigurename}
\listoffigures
% (El "Índice de gráficas" de la guía solo procede si separamos las gráficas
% de las ilustraciones con un tipo de float propio; de momento van juntas.)
% =========================================================================
% Lista de acrónimos (paquete acronym).
% Definir cada acrónimo aquí y usarlo en el texto con \ac{CLAVE}: la
% primera vez imprime la forma completa ("Trabajo de Fin de Grado (TFG)")
% y las siguientes solo la sigla ("TFG"). Otras variantes: \acl (largo),
% \acs (sigla), \acp (plural).
% =========================================================================

\chapter*{Listado de acrónimos}
\addcontentsline{toc}{chapter}{Listado de acrónimos}

\begin{acronym}[BUKIT]  % el argumento fija el ancho de la columna de siglas
  \acro{API}{\textit{Application Programming Interface}}
  \acro{ODS}{Objetivos de Desarrollo Sostenible}
  \acro{TFG}{Trabajo de Fin de Grado}
\end{acronym}
             % Listado de acrónimos
% =========================================================================
% Empleo de Inteligencia Artificial (IA) — apartado preliminar V de la
% guía definitiva de la EIB (si procede; en este TFG procede claramente).
% Declarar qué herramientas de IA se han empleado y para qué: asistencia
% en programación (p. ej. Claude Code), redacción, generación de
% diagramas... Transparencia total: es un requisito del centro.
% =========================================================================

\chapter*{Empleo de Inteligencia Artificial (IA)}
\addcontentsline{toc}{chapter}{Empleo de Inteligencia Artificial (IA)}

\todonote{Contenido pendiente: declarar las herramientas de IA empleadas
y su propósito en el desarrollo y la redacción.}
            % Empleo de IA (guía EIB, apdo. V)

\tableofcontents                         % Índice de la memoria

% ----- MEMORIA (numeración arábiga desde aquí) -----
\clearpage
\pagenumbering{arabic}

% =========================================================================
% Capítulo 1 — Introducción
% Según la guía del TFG, la introducción debe exponer:
%   - los objetivos del TFG (ámbito, alcance y límites; aquí solo se
%     presentan: se desarrollan en su propio capítulo),
%   - el estado de la cuestión y las circunstancias de la elección,
%   - la pertinencia del tema y los motivos de su elección,
%   - la metodología empleada,
%   - la enumeración y descripción general de los capítulos, con su
%     justificación desde el punto de vista de la estructura.
% =========================================================================

\chapter{Introducción}
\label{ch:introduccion}

\section{El problema y la propuesta}
\label{sec:intro-problema}

Casi todo el mundo mantiene una lista, escrita o mental, de cosas que quiere hacer en la vida: viajes, hábitos, proyectos, experiencias —las llamadas \eng{bucket lists} o listas de deseos vitales—.
La mayoría de esas metas se queda en deseo: la intención por sí sola explica solo una parte pequeña de la conducta \parencite{sheeran-intention}, y las condiciones que la convierten en acción —plazos concretos, seguimiento del progreso, apoyo de otras personas— rara vez acompañan a una lista guardada en un cajón.
Las redes sociales dominantes no ayudan: su experiencia está volcada en el consumo de contenido, no en la acción —una lectura que el capítulo \ref{ch:beneficios} desarrolla y fundamenta—, y su uso pasivo está asociado a peor bienestar \parencite{verduyn-sns}.

\bukit es la respuesta que este \ac{TFG} construye: una red social cuyo propósito es ayudar a cumplir metas vitales y a no perseguirlas en solitario.
Toda la experiencia se orienta a la acción sobre tres pilares: estructura (metas agrupadas en listas con plazo y recordatorios), comunidad (personas con objetivos afines que se acompañan) y recompensa al avance (visibilidad para quien progresa, no para quien más audiencia acumula).
El capítulo \ref{ch:contexto} desarrolla el concepto completo y el \ref{ch:beneficios} fundamenta con literatura por qué ese diseño puede funcionar.

El resultado material del trabajo es una aplicación real: una \ac{API} propia con Node y PostgreSQL y un cliente web con mapa social, mensajería y verificación de identidad —el compromiso mínimo del trabajo—, acompañados como objetivos secundarios de una aplicación móvil parcial y de funciones asistidas por \ac{IA}, con la lógica compartida entre clientes desde paquetes comunes.
Los objetivos comprometidos, con su alcance y sus límites exactos, se definen en el capítulo \ref{ch:objetivos}.

\section{Pertinencia del tema y motivos de la elección}
\label{sec:intro-motivos}

Los motivos de la elección son tres.
El primero es la continuidad: el autor ya construyó y documentó TrackBites \parencite{deleanu-trackbites}, una aplicación web completa de registro nutricional, y este trabajo da el salto de dificultad que aquella experiencia deja preparado —de una aplicación web en JavaScript a un producto multiplataforma en TypeScript con base de datos relacional, geolocalización, búsqueda por afinidad y despliegue continuo (capítulo \ref{ch:alternativas})—.
El segundo es la cobertura de las competencias del grado: un producto de este tipo obliga a ejercitar en un solo proyecto el diseño de datos, la construcción de una \ac{API} segura, dos clientes con lógica compartida, las pruebas automatizadas, la entrega continua y la gestión completa del proyecto (planificación, riesgos, seguimiento y presupuesto) que esta memoria documenta.
El tercero es personal y el capítulo \ref{ch:beneficios} lo declara abiertamente: el autor quiere una herramienta que empuje a las personas de consumir a crear, y ese propósito existía antes que este \ac{TFG}.

\section{Metodología}
\label{sec:intro-metodologia}

El proyecto se gestiona con un método híbrido de la familia Scrumban adaptado a una sola persona: iteraciones semanales con capacidad variable declarada, tablero con límite de trabajo en curso, historias de usuario estimadas en puntos y medición continua de la velocidad para recalibrar el plan en cada cierre (capítulo \ref{ch:tareas}).
El desarrollo se realiza con asistencia declarada de \ac{IA} generativa, con salvaguardas explícitas contra la dependencia acrítica (riesgo R5, capítulo \ref{ch:riesgos}); cada decisión queda registrada en un diario de desarrollo con su alternativa descartada y su porqué, y esa trazabilidad es la materia prima de esta memoria.
La gestión de riesgos es un proceso continuo con seguimiento semanal (capítulo \ref{ch:riesgos}) y el gasto real del trabajo se registra mes a mes (capítulo \ref{ch:presupuesto}).

\section{Estructura de la memoria}
\label{sec:intro-estructura}

La memoria sigue el índice oficial de la \ac{EIB} y su lógica interna va de la justificación (qué se hace y por qué) a la solución (cómo), de ahí a la gestión (cuándo y a qué coste) y cierra con la evidencia de lo logrado.
El bloque de análisis lo forman el contexto y el concepto del producto (capítulo \ref{ch:contexto}), los objetivos y el alcance (capítulo \ref{ch:objetivos}), los beneficios y su relación con los \ac{ODS} (capítulo \ref{ch:beneficios}), el estado del arte con los requisitos del sistema (capítulo \ref{ch:estado-del-arte}), el análisis de alternativas tecnológicas (capítulo \ref{ch:alternativas}) y el plan de riesgos (capítulo \ref{ch:riesgos}).
El bloque de solución describe la propuesta (capítulo \ref{ch:solucion}) y su diseño de alto nivel (capítulo \ref{ch:diseno}).
El bloque de gestión y ejecución recoge la metodología con la planificación y la estimación del esfuerzo (capítulo \ref{ch:tareas}), los cálculos y algoritmos del producto (capítulo \ref{ch:algoritmos}), el análisis de resultados (capítulo \ref{ch:resultados}), el cronograma completo del \ac{TFG} (capítulo \ref{ch:gantt-tfg}) y el presupuesto (capítulo \ref{ch:presupuesto}).
Cierran las conclusiones y el trabajo futuro (capítulo \ref{ch:conclusiones}).
         % obligatorio
% =========================================================================
% Capítulo 2 — Contexto (OBLIGATORIO)
% Situación en la que nace el proyecto: el problema de las redes sociales
% de entretenimiento vs. una red orientada a cumplir objetivos, el marco
% académico (TFG) y personal, etc.
% (La relación con los ODS va en el capítulo 4, junto con los beneficios,
% según la guía definitiva de la EIB.)
% =========================================================================

\chapter{Contexto}
\label{ch:contexto}

\section{El problema: metas que se abandonan y redes que no acompañan}
\label{sec:contexto-problema}

Las metas vitales se abandonan siguiendo un patrón reconocible: deseos formulados sin plazo ni medida, progreso que nadie registra y un camino que se recorre en solitario, sin nadie alrededor que acompañe o pida cuentas.
El capítulo \ref{ch:beneficios} fundamenta esa lectura con la literatura de la psicología de la motivación; este capítulo la da por presentada para situar el producto.

Las redes sociales de entretenimiento no alivian el problema y hay razones fundadas —presentadas en el capítulo \ref{ch:beneficios}— para pensar que lo agravan: su experiencia gira en torno al desplazamiento infinito por contenido ajeno y a la comparación con vidas editadas.
El hueco que \bukit ocupa es precisamente ese: entre los productos que el capítulo \ref{ch:estado-del-arte} analiza no hay una red social generalista cuya experiencia completa esté orientada a que la persona \emph{haga}, en lugar de a que mire.

\section{La propuesta: una red social orientada a la acción}
\label{sec:contexto-propuesta}

\bukit se organiza en cuatro pestañas, el patrón de navegación asentado por las redes dominantes \parencite{instagram-app}, para que la barrera de entrada sea la de una aplicación ya conocida:

\begin{enumerate}
  \item \textbf{Inicio}: un \eng{feed} —el flujo de publicaciones de la pantalla de inicio— personalizado en mosaico con ideas de metas, progresos, metas completadas y no logradas y vídeos de motivación, más un buscador de perfiles, publicaciones y listas; deslizando, un segundo \eng{feed} solo con el contenido de las personas seguidas.
  \item \textbf{Listas}: el centro de control de los objetivos: crear, ver, editar y eliminar listas de metas con plazo, con recordatorios que envía la propia aplicación.
  \item \textbf{Mensajes y mapa}: la mensajería, con un mini-mapa en la parte superior; el mapa a pantalla completa ofrece dos vistas —la de seguidos mutuos y la de descubrimiento de personas afines, filtrable por país, por radio de cercanía o en todo el mundo— y aquí viven también las notificaciones.
  \item \textbf{Perfil}: quién eres: hasta cinco publicaciones de tema libre —las únicas no ligadas a listas— e historias ilimitadas.
\end{enumerate}

Sobre esa estructura, cinco reglas de producto concretan los tres pilares del concepto —estructura, comunidad y recompensa al avance— y condicionan el diseño técnico del resto de la memoria:

\begin{itemize}
  \item \textbf{Las listas son la unidad de las metas.}
        Una lista agrupa metas con plazo, y las publicaciones de progreso o de meta completada solo pueden crearse desde listas propias: nadie puede presumir de un avance que no persigue de verdad.
        El plan gratuito admite un máximo de tres listas.
  \item \textbf{La visibilidad premia el avance.}
        Hay dos tipos de usuario: normal (cualquier persona registrada) y TOP (creador admitido por formulario externo, con vídeo ilimitado y mayor peso en el \eng{feed}); la ordenación del \eng{feed} pondera el progreso y la aportación, no la mera popularidad, y los vídeos de motivación son exclusivos de los perfiles TOP.
  \item \textbf{El contenido de perfil es efímero.}
        Todo lo publicado en el perfil caduca a los 60 días salvo que su autor pulse «mantener»: el escaparate personal no es el centro de la experiencia.
  \item \textbf{Conocer gente exige verificación.}
        La verificación de identidad —documento de identidad, teléfono y fotografías de validación, delegada en un proveedor especializado— desbloquea publicar desde el perfil y el mapa de descubrimiento, reservado a personas mayores de edad, con ubicación siempre aproximada y con consentimiento explícito.
  \item \textbf{Los recordatorios empujan.}
        La aplicación recuerda los plazos de las metas: el compromiso es medible y con fecha, no un deseo vago.
\end{itemize}

La traducción de estas reglas a beneficios esperados, con su fundamento en la literatura, está en el capítulo \ref{ch:beneficios}; su traducción a requisitos del sistema, en el capítulo \ref{ch:estado-del-arte}.

\section{Marco académico y personal}
\label{sec:contexto-marco}

Este proyecto es el \ac{TFG} del Grado en Ingeniería Informática de Gestión y Sistemas de Información de la \ac{EIB}, centro de la \ac{UPVEHU}, con 12 créditos \ac{ECTS} —unas 360 horas de referencia académica, cuya conciliación con la dedicación realmente planificada se recoge en el capítulo \ref{ch:gantt-tfg}— y desarrollo previsto a lo largo del curso 2026--2027 completo.
El calendario está condicionado por circunstancias personales declaradas en la planificación: compromisos laborales durante el verano, una estancia Erasmus en el primer cuatrimestre y prácticas de empresa desde febrero (capítulos \ref{ch:tareas} y \ref{ch:gantt-tfg}).

El trabajo no parte de cero.
Su precedente directo es TrackBites \parencite{deleanu-trackbites}, la aplicación web completa que el autor construyó con anterioridad y cuya arquitectura funciona como línea base de las decisiones tecnológicas de este proyecto (capítulo \ref{ch:alternativas}).
El desarrollo se realiza en solitario y con la asistencia de \ac{IA} generativa ya acotada en la metodología (apartado \ref{sec:intro-metodologia}).
             % obligatorio
% =========================================================================
% Capítulo 3 — Objetivos y alcance del trabajo (OBLIGATORIO)
% Objetivo principal + objetivos secundarios (medibles), qué entra y qué
% queda explícitamente FUERA del alcance (límites), y criterios de éxito.
% =========================================================================

\chapter{Objetivos y alcance del trabajo}
\label{ch:objetivos}

Los objetivos de este \ac{TFG} se formulan siguiendo el marco \ac{SMART}, es decir, como compromisos específicos, medibles, alcanzables, relevantes y acotados en el tiempo \parencite{doran-smart}.
Se opta por este marco frente a una enumeración informal de intenciones porque convierte cada objetivo en un criterio verificable: al cierre del trabajo puede comprobarse de forma objetiva si se ha cumplido o no.

El horizonte temporal común a todos los objetivos es la defensa del trabajo en la convocatoria de junio de \thesisYear.
El desarrollo se concentra en un periodo de ejecución de unos tres meses, con una dedicación de referencia de 360 horas, correspondiente a los 12 créditos \ac{ECTS} de la asignatura.
La planificación detallada de tareas, fechas e hitos se recoge en los capítulos \ref{ch:tareas} y \ref{ch:gantt-tfg}.

Los objetivos se organizan en dos bloques: los \emph{técnicos y de producto}, que describen qué debe existir y funcionar al terminar, y los \emph{de gestión y académicos}, que describen cómo debe ejecutarse y documentarse el trabajo.
Dentro del primer bloque se distingue entre objetivos comprometidos y objetivos secundarios: los segundos aportan valor y validan decisiones de diseño, pero su ausencia no compromete el éxito del trabajo.

\section{Objetivo principal}
\label{sec:objetivo-principal}

Diseñar, desarrollar y desplegar \bukit, una red social orientada a ayudar a cumplir objetivos vitales (\eng{bucket lists}), entregando antes de la defensa un \eng{backend} completo (\ac{API} \ac{REST} y base de datos) y un cliente web funcional que cubra la experiencia de producto definida, junto con una memoria que documente y justifique todas las decisiones tomadas.

\section{Objetivos técnicos y de producto}
\label{sec:objetivos-tecnicos}

\begin{itemize}
  \item \textbf{Backend único para las tres necesidades del producto.}
        Implementar con Express y Prisma una \ac{API} \ac{REST} sobre una única base de datos PostgreSQL que resuelva el grafo social (integridad relacional), el mapa (consultas geoespaciales por radio con PostGIS) y la afinidad entre personas usuarias (búsqueda por similitud vectorial con pgvector).
        El objetivo se considera cumplido cuando el 100\,\% de los flujos núcleo del producto (cuentas, listas y metas, publicaciones, seguimiento, mensajería, mapa y notificaciones) está expuesto por la \ac{API} y cubierto por pruebas de integración antes del cierre del periodo de ejecución.
  \item \textbf{Autenticación y seguridad desde el diseño.}
        Proteger la \ac{API} con autenticación \ac{JWT} con rotación de \eng{refresh tokens}, contraseñas cifradas con Bcrypt, doble capa de limitación de peticiones y manejo centralizado de errores.
        Medible: el 100\,\% de los \eng{endpoints} privados exige un token válido, y los casos de abuso básicos (credenciales inválidas, exceso de peticiones) quedan cubiertos por pruebas automatizadas.
  \item \textbf{Cliente web completo.}
        Desarrollar con React y Vite una \ac{SPA} que implemente las cuatro pestañas del producto: inicio con \eng{feed} personalizado, listas, mensajes con mapa y notificaciones, y perfil.
        Medible: una persona usuaria puede completar de principio a fin los flujos núcleo desde el navegador, sin errores bloqueantes, antes del cierre del periodo de ejecución.
  \item \textbf{Lógica compartida sin duplicación.}
        Centralizar los tipos, la validación (Zod) y el cliente HTTP en los paquetes compartidos del \eng{monorepo} (\texttt{packages/core}, \texttt{packages/api-client}, \texttt{packages/ui-core}) como única fuente de verdad para web y móvil.
        Medible: cero esquemas de validación duplicados entre aplicaciones al cierre del desarrollo.
  \item \textbf{Aplicación móvil parcial (objetivo secundario).}
        Construir con React Native y Expo las pantallas núcleo (autenticación, listas y \eng{feed}) reutilizando la lógica compartida, como demostración de la arquitectura multiplataforma.
        Se declara secundario porque el compromiso mínimo del trabajo es \eng{backend} más web; completarlo valida la decisión arquitectónica central, pero no completarlo no invalida el trabajo.
  \item \textbf{Estimación de plazos asistida por \ac{IA} (objetivo secundario).}
        Integrar la estimación de plazos realistas por meta mediante un modelo de lenguaje con salida estructurada, ejecutada en segundo plano y dirigida por eventos (nunca por sondeo).
        Se declara secundario porque depende de servicios externos cuyo coste y disponibilidad quedan fuera del control del proyecto.
  \item \textbf{Calidad y despliegue continuo.}
        Montar una cadena de integración continua con estrategia de fallo rápido (análisis estático, comprobación de tipos, pruebas unitarias y de integración, y pruebas \ac{E2E} como última barrera) con despliegue automático al superar todas las etapas.
        Medible: la rama principal se mantiene en verde y cada cambio validado llega a producción sin intervención manual.
\end{itemize}

\section{Objetivos de gestión y académicos}
\label{sec:objetivos-gestion}

\begin{itemize}
  \item \textbf{Dedicación y plazo.}
        Ejecutar el desarrollo en un periodo de unos tres meses, con una dedicación total en el entorno de las 360 horas correspondientes a los 12 créditos \ac{ECTS}, finalizando producto y memoria antes de la convocatoria de junio de \thesisYear.
        Medible: el registro de horas por sesión del diario de desarrollo del repositorio.
  \item \textbf{Trazabilidad de las decisiones.}
        Registrar el 100\,\% de las sesiones de trabajo en el diario de desarrollo, anotando fecha, horas invertidas, problemas encontrados y cada decisión junto con la alternativa descartada y su porqué.
        Este diario es la fuente de datos real de la planificación (capítulos \ref{ch:tareas} y \ref{ch:gantt-tfg}), por lo que su exhaustividad condiciona directamente la calidad de la memoria.
  \item \textbf{Memoria en paralelo al desarrollo.}
        Redactar la memoria de forma incremental, junto al desarrollo y no al final, manteniendo el documento siempre compilable sin errores.
        Medible: al cierre de cada fase del desarrollo, los capítulos que la documentan contienen redacción real y no marcadores de contenido pendiente.
  \item \textbf{Empleo documentado de la \ac{IA} generativa.}
        Emplear herramientas de \ac{IA} generativa como asistentes del desarrollo y declarar su uso conforme a la guía de la \ac{EIB}, analizando de forma crítica su aportación real, sus límites y los problemas encontrados.
        Medible: la declaración de empleo de \ac{IA} de esta memoria enumera todas las herramientas utilizadas y los usos que se les han dado.
\end{itemize}

\section{Alcance}
\label{sec:alcance}

Delimitar el alcance es tan importante como enumerar los objetivos: fija qué se compromete, asegura la entrega de un producto mínimo funcional y protege el trabajo frente a la expansión descontrolada de requisitos (\eng{feature creep}).

\subsection{Dominio de la aplicación}
\label{sec:alcance-dominio}

El sistema se contextualiza en el dominio de una red social orientada a la acción: ayudar a cumplir objetivos vitales (\eng{bucket lists}) y a no perseguirlos en solitario.
Los elementos principales sobre los que orbita el sistema son las personas usuarias, las listas de metas con plazo, las publicaciones de progreso y las conexiones entre personas afines (seguimiento, mensajería y mapa).
La experiencia se organiza en cuatro pestañas (inicio, listas, mensajes y perfil), y esa misma estructura ordena los módulos funcionales del alcance.

\subsection{Requisitos funcionales (dentro del alcance)}
\label{sec:alcance-dentro}

Para garantizar la viabilidad del trabajo, el desarrollo se limita a los siguientes módulos:

\begin{itemize}
  \item \textbf{Usuarios, roles y verificación.}
        Registro e identificación única, con dos tipos de usuario (normal y TOP) que reciben distinta ponderación en el \eng{feed}.
        La verificación de identidad se delega en un proveedor externo en modo de pruebas, almacenando únicamente el estado y el \eng{token} de la verificación por minimización de datos conforme al \ac{RGPD}; superarla desbloquea publicar desde el perfil y el mapa de descubrimiento, restringido a personas mayores de edad.
  \item \textbf{Listas y metas (unidad central del producto).}
        Creación, consulta, edición y eliminación de listas de metas con plazo, con recordatorios asociados y un máximo de tres listas en el plan gratuito.
  \item \textbf{Publicaciones y \eng{feed}.}
        Publicaciones de progreso y de meta completada, creadas únicamente desde listas propias; \eng{feed} personalizado en mosaico con ponderación por tipo de usuario; segundo \eng{feed} solo de seguidos; y buscador de perfiles, publicaciones y listas.
  \item \textbf{Mensajería, mapa y notificaciones.}
        Chat en tiempo real entre personas usuarias; mapa con vista de seguidos mutuos y vista de descubrimiento por proximidad, apoyada en el cálculo de afinidad basado en \eng{embeddings}; y centro de notificaciones.
  \item \textbf{Perfil.}
        Hasta cinco publicaciones libres e historias ilimitadas, con contenido efímero que se elimina a los 60 días salvo acción expresa de mantenerlo.
  \item \textbf{Infraestructura transversal.}
        El \eng{monorepo} completo, con la aplicación móvil en su versión parcial (objetivo secundario, sección \ref{sec:objetivos-tecnicos}) y los paquetes compartidos de lógica; la estimación de plazos asistida por \ac{IA}, dirigida por eventos; y la infraestructura de calidad: pruebas unitarias, de integración y \ac{E2E}, con integración y despliegue continuos.
\end{itemize}

\subsection{Funcionalidades excluidas (fuera del alcance)}
\label{sec:alcance-fuera}

Para acotar el esfuerzo y proteger los plazos, quedan explícitamente fuera del alcance de este \ac{TFG}, como líneas futuras:

\begin{itemize}
  \item \textbf{Planes de pago.}
        No se implementan pasarelas de pago ni lógica de suscripción; el límite de tres listas del plan gratuito queda fijo.
  \item \textbf{Verificación de identidad en producción.}
        El proveedor externo se integra únicamente en modo de pruebas; su contratación real y los costes asociados quedan fuera.
  \item \textbf{Red de anuncios real.}
        No se integra ninguna red comercial; el hueco publicitario del \eng{feed} se cubre con una tabla propia de publicaciones patrocinadas.
  \item \textbf{Soporte de vídeo.}
        No se implementan transcodificación ni \eng{streaming}; el contenido multimedia se limita a imágenes.
  \item \textbf{Publicación en tiendas de aplicaciones.}
        La distribución en App Store y Play Store queda fuera; la aplicación móvil se valida con compilaciones de desarrollo.
  \item \textbf{Funciones aún en definición.}
        La penalización social por metas incumplidas y el panel de administración de usuarios TOP se posponen hasta que su diseño madure.
\end{itemize}

Estas líneas se retoman como trabajo futuro en el capítulo \ref{ch:conclusiones}.

\section{Criterios de éxito}
\label{sec:criterios-exito}

El trabajo se considerará un éxito si, a su cierre, se cumplen de manera binaria (sí/no) las siguientes condiciones:

\begin{enumerate}
  \item \textbf{Producto funcional desplegado.}
        Una persona usuaria puede registrarse, crear una lista con metas y plazo, publicar un progreso desde su lista, seguir a otras personas, consultar su \eng{feed}, conversar por chat y usar el mapa, sin errores bloqueantes y sobre el despliegue de producción.
  \item \textbf{Arquitectura compartida demostrada.}
        Los tipos y la validación viven una única vez en los paquetes compartidos y ninguna aplicación duplica esquemas.
  \item \textbf{Calidad automatizada.}
        La cadena de integración y despliegue continuos está operativa, en verde sobre la rama principal, y cada cambio validado se despliega sin intervención manual.
  \item \textbf{Trazabilidad completa.}
        El diario de desarrollo cubre todas las sesiones de trabajo y la memoria completa el índice oficial de la \ac{EIB} sin contenidos pendientes, con las decisiones justificadas frente a sus alternativas.
  \item \textbf{Ejecución dentro del marco académico.}
        El desarrollo se ha completado dentro del periodo planificado y la dedicación registrada es coherente con las 360 horas de los 12 créditos \ac{ECTS}.
\end{enumerate}
    % obligatorio
% =========================================================================
% Capítulo 4 — Beneficios que aporta el trabajo y su relación con los ODS
% (OBLIGATORIO)
% Beneficios técnicos, sociales y económicos, y su conexión con los
% Objetivos de Desarrollo Sostenible (p. ej. salud y bienestar, reducción
% de desigualdades), justificándola.
% =========================================================================

\chapter{Beneficios que aporta el trabajo y su relación con los Objetivos
  de Desarrollo Sostenible}
\label{ch:beneficios}

\todonote{Contenido pendiente.}

\section{Relación con los Objetivos de Desarrollo Sostenible}
\label{sec:ods}

\todonote{Contenido pendiente.}
       % obligatorio
% =========================================================================
% Capítulo 5 — Descripción de requerimientos y análisis del estado del arte
% (si procede — para BUKIT procede)
% Estado del arte: redes sociales existentes (Instagram, Pinterest, apps
% de hábitos/objetivos) y qué hueco cubre BUKIT. Requerimientos:
% funcionales y no funcionales del sistema.
% =========================================================================

\chapter{Descripción de requerimientos y análisis del estado del arte}
\label{ch:estado-del-arte}

\todonote{Contenido pendiente.}
      % si procede (procede)
% =========================================================================
% Capítulo 6 — Análisis de alternativas
% Método: para cada ámbito, qué se necesita, qué opciones había, cuál se
% elige y por qué. El punto de partida es siempre la arquitectura de
% TrackBites (proyecto previo del autor); las decisiones aún abiertas se
% declaran como tales, con sus candidatas y el criterio que decidirá.
% Fuentes: propuesta original (§4), dev-journal, CLAUDE.md y el propio
% estado del repositorio. La solución resultante se describe en el cap. 8.
% =========================================================================

\chapter{Análisis de alternativas}
\label{ch:alternativas}

Este capítulo recorre las decisiones tecnológicas del proyecto: para cada ámbito se expone qué se necesita, qué opciones se han considerado, cuál se elige y por qué.
El análisis no parte de cero.
El autor desarrolló previamente TrackBites \parencite{deleanu-trackbites} (apartado \ref{sec:contexto-marco}), y aquella arquitectura funciona aquí como línea base: cada elección se plantea como «lo ya dominado» frente a lo que este producto exige, y solo se salta donde el salto aporta.
La propuesta de este \ac{TFG} concentró ese salto en cuatro mejoras: TypeScript en todo el código, lógica compartida entre web y móvil, una base de datos relacional con capacidades geoespaciales y vectoriales, y el cierre del ciclo de entrega con despliegue automático.

Tres criterios transversales, fijados en capítulos anteriores, atraviesan todas las decisiones.
Primero, el \ac{TFG} debe demostrar las competencias del grado, así que las piezas centrales se construyen a mano en lugar de delegarse en plataformas que las regalan hechas (sección \ref{sec:alt-supabase}).
Segundo, la restricción de coste cero: solo se adoptan servicios con capa gratuita suficiente para un uso académico (riesgo R9, capítulo \ref{ch:riesgos}).
Tercero, el \ac{RGPD} condiciona dónde viven los datos (región europea) y qué servicios pueden tratarlos (riesgo R7).
Las decisiones diferidas se recogen con sus candidatas y con el criterio que las cierra.
La solución resultante de este análisis se describe en el capítulo \ref{ch:solucion}.

\section{Lenguaje: TypeScript en todo el código}
\label{sec:alt-lenguaje}

TrackBites se escribió en JavaScript.
Aquí todo el código —servidor, web, móvil y paquetes compartidos— pasa a TypeScript, la decisión que más mantenibilidad aporta a medio y largo plazo: el compilador detecta antes de ejecutar errores que en JavaScript solo aparecen en ejecución \parencite{typescript-docs}, y los tipos son la base técnica de la validación compartida de la sección \ref{sec:alt-validacion}.
La configuración es deliberadamente estricta (modo \eng{strict} con comprobaciones adicionales) y única para todo el repositorio, de modo que ningún paquete pueda relajarla por su cuenta.

\section{Estructura del repositorio: \eng{monorepo} con pnpm}
\label{sec:alt-monorepo}

Compartir tipos, validación y cliente de \ac{API} entre tres aplicaciones exige que todo el código viva junto: un único repositorio con las aplicaciones (\texttt{apps/}) y los paquetes compartidos (\texttt{packages/}), el llamado \eng{monorepo}.
La alternativa de un repositorio por aplicación se descarta porque cada pieza compartida tendría que publicarse y versionarse para poder consumirse, un peaje continuo sin sentido para una sola persona.

Dentro del \eng{monorepo} compitieron dos herramientas.
La propuesta original apuntaba a Turborepo, un orquestador que acelera \eng{monorepos} grandes con caché de tareas y compilaciones incrementales \parencite{turborepo-docs}.
Al montar el repositorio se optó en su lugar por los \eng{workspaces} de pnpm sin orquestador \parencite{pnpm-docs}.
La arquitectura elegida explica por sí sola esa renuncia: los paquetes compartidos se consumen \emph{sin compilar} (sus puntos de entrada apuntan directamente al código fuente y cada aplicación los compila con su propia herramienta), así que apenas quedan tareas que orquestar ni compilaciones que cachear, y la capa extra de Turborepo pagaría complejidad sin devolver velocidad en un equipo de una persona.
pnpm aporta además la reproducibilidad que el proyecto se exige: versiones de dependencias fijadas exactas (sin rangos) con el archivo de bloqueo versionado en el repositorio, y versiones de Node y del propio gestor clavadas en la configuración, de modo que cualquier máquina reconstruye el mismo entorno.
Queda preparada una contingencia documentada: si el empaquetador de React Native no tolerase los enlaces simbólicos con los que pnpm instala, una línea comentada en la configuración cambia a una instalación plana equivalente a la de npm sin perder el almacén rápido de pnpm.

\section{Base de datos: PostgreSQL con PostGIS y pgvector}
\label{sec:alt-bd}

Es la decisión más relevante del conjunto, y la única que hereda una lección previa documentada.
TrackBites usó MongoDB con Mongoose, y su memoria dejó registrado un compromiso consciente: aquel dominio, estrictamente jerárquico, encajaba mejor en una base relacional, y la integridad referencial y los borrados en cascada hubo que emularlos a mano con ganchos del propio Mongoose.
\bukit lleva esa conclusión a la práctica.
Su dominio es un grafo social —personas que se siguen, listas, metas, publicaciones ligadas a listas— con dos necesidades especiales: buscar personas por cercanía geográfica y por afinidad de intereses.
Un modelo relacional resuelve el grafo con claves foráneas e integridad nativas, y aprovecha la base sólida en \ac{SQL} del autor.

Las dos necesidades especiales no obligan a añadir motores dedicados, porque PostgreSQL las cubre con extensiones: PostGIS añade tipos y consultas geoespaciales —búsqueda por radio con índice espacial— \parencite{postgis-docs}, y pgvector añade columnas vectoriales y búsqueda por similitud \parencite{pgvector-docs} sobre las representaciones numéricas (\eng{embeddings}) de metas e intereses con las que se calcula la afinidad.
La alternativa de sumar servicios especializados (una base vectorial y un servicio geoespacial junto a la base principal) multiplicaría las piezas que operar, pagar y sincronizar; un único PostgreSQL cubre las tres necesidades a la vez.

Para producción se elige Neon, un PostgreSQL gestionado \parencite{neon-docs}, entre las candidatas de la propuesta (Neon, Supabase, Railway): la residencia europea del dato es el criterio registrado de esa elección, y su capa gratuita satisface además la restricción de coste cero.
La disposición operativa de la base —Docker en desarrollo, pruebas e integración continua; Neon en producción— se describe en el capítulo \ref{ch:solucion}.

\section{Acceso a datos: Prisma}
\label{sec:alt-orm}

Como capa de acceso a datos (\ac{ORM}) se adopta Prisma, que aporta migraciones versionadas del esquema y generación automática de los tipos TypeScript de cada entidad \parencite{prisma-docs}: cada cambio del modelo queda versionado junto al código, y los tipos del servidor no pueden desincronizarse de la base.
Sustituye a Mongoose, la capa equivalente del mundo MongoDB, que queda descartada junto con su motor.
Prisma no cubre las consultas geoespaciales ni vectoriales de las extensiones; esas consultas concretas se escriben en \ac{SQL} crudo desde el propio Prisma, una convivencia prevista desde el diseño.
Las extensiones se declaran además en el propio esquema de Prisma, de modo que el código y la base nunca discrepen sobre qué hay activado.

\section{Servidor de la \acs{API}: Express 5 frente a NestJS}
\label{sec:alt-servidor}

La propuesta dejó abierta la elección entre NestJS y «Express bien estructurado», y el arranque del repositorio la cerró a favor de Express en su versión 5.
NestJS ofrece una arquitectura completa y opinada —módulos, inyección de dependencias, decoradores— \parencite{nestjs-docs}; Express es minimalista y deja la estructura en manos del proyecto \parencite{express-docs}.
Esa diferencia basta para sostener la elección.
El valor académico de este \eng{backend} está justamente en lo hecho a mano —autenticación con \ac{JWT} y rotación de los \eng{tokens} de refresco, doble capa de limitación de peticiones, manejo de errores centralizado: patrones que el autor ya domina de TrackBites, precisamente sobre Express—, y NestJS los entregaría envueltos en las abstracciones del marco, restando demostración y sumando su propia curva de aprendizaje a la del resto de la pila (riesgo R6, capítulo \ref{ch:riesgos}).
\section{Pruebas de la \acs{API} y registro de sucesos}
\label{sec:alt-pruebas-registro}

Dos decisiones de menor calado acompañan al servidor, ambas con alternativa heredada de TrackBites o de su ecosistema.
Las pruebas de integración de la \ac{API} se escriben con Vitest \parencite{vitest-docs} y Supertest, frente al precedente de TrackBites (Jest con Supertest).
Dos razones cierran la elección: un único ejecutor de pruebas para todo el \eng{monorepo} (los paquetes compartidos ya prueban con Vitest) y el hecho de que la \ac{API} se escribe como módulos ECMAScript, un terreno donde Jest exige una configuración frágil.
El registro de sucesos se resuelve con pino \parencite{pino-docs} y su \eng{middleware} pino-http, frente a winston: pino produce de forma nativa registro estructurado en \ac{JSON} con niveles de gravedad, su \eng{middleware} \ac{HTTP} asigna un identificador a cada petición para seguirla en los registros, y su rendimiento es mejor.

\section{La alternativa global: Supabase}
\label{sec:alt-supabase}

Antes de elegir pieza a pieza cabía una enmienda a la totalidad.
Supabase empaqueta en un solo producto casi todo lo que este \eng{backend} construye: PostgreSQL gestionado con PostGIS y pgvector, autenticación, almacenamiento de archivos, tiempo real y funciones en la nube \parencite{supabase-docs}.
Su ventaja es real y se reconoce: desarrollo mucho más rápido y menos infraestructura que mantener.
Se descarta por la razón inversa: menos \eng{backend} hecho a mano, que es justo lo que el autor domina y quiere demostrar.
Un \ac{TFG} cuyo \eng{backend} fuese configuración de Supabase acreditaría integración, no diseño; el \eng{backend} propio con Node, Prisma y PostgreSQL es el núcleo que demuestra las competencias del grado.
Supabase queda anotada como vía de aceleración legítima para una eventual segunda versión, ya fuera de este \ac{TFG}.

\section{Validación compartida: Zod frente a Joi}
\label{sec:alt-validacion}

TrackBites validaba con Joi solo en el servidor, sin validación compartida con el cliente.
Aquí la validación vive una única vez, en el paquete compartido \texttt{packages/core}, escrita con Zod: cada esquema valida los datos en ejecución y de él se infiere automáticamente el tipo TypeScript \parencite{zod-docs}, de modo que servidor, web y móvil comparten literalmente la misma regla y el mismo tipo.
Con Joi, validación y tipos habrían tenido que mantenerse por separado \parencite{joi-docs}: el doble mantenimiento que este diseño elimina es exactamente el porqué del cambio.
La regla de proyecto lo fija: ningún esquema se duplica en las aplicaciones.

\section{Cliente web: React y Vite, con TanStack Query}
\label{sec:alt-web}

La base del cliente web se conserva de TrackBites: React con Vite, ahora en TypeScript.
Cambiar de marco (por ejemplo, a uno con renderizado en el servidor) no aportaría a una aplicación que vive tras un inicio de sesión, y desperdiciaría la experiencia acumulada.
El cambio está en el estado de servidor.
En TrackBites la caché de datos remotos se gestionaba a mano con Redux Toolkit, podando el árbol de estado local para evitar recargas; funcionó, pero a costa de mucho código repetitivo de carga, error y caché.
Para un producto cuyo centro es un \eng{feed} con desplazamiento infinito, esa gestión manual se descarta en favor de TanStack Query, que resuelve de serie la caché, los reintentos, la invalidación y la paginación infinita del estado de servidor \parencite{tanstack-query-docs}.
La biblioteca de estado global de interfaz (Redux Toolkit con Reselect, en consideración por la experiencia previa) se decide al montar el cliente web, con el primer estado global real como caso de prueba.
Por el mismo principio, el andamiaje de las aplicaciones web y móvil se genera con sus herramientas oficiales en el momento de abordarlas, de modo que la configuración nace alineada con la versión vigente de cada herramienta.

\section{Cliente móvil: React Native con Expo}
\label{sec:alt-movil}

La aplicación móvil es objetivo secundario (capítulo \ref{ch:objetivos}), pero su elección condiciona la arquitectura desde el primer día.
Se elige React Native con Expo.
React Native permite escribir la interfaz móvil en el mismo lenguaje y paradigma que la web \parencite{react-native-docs}, con lo que los paquetes compartidos (tipos, validación, cliente de \ac{API}, lógica de estado) sirven a ambas plataformas sin traducción; el desarrollo nativo por plataforma —o un marco multiplataforma de otro lenguaje— rompería exactamente esa compartición, que es el principio rector de la arquitectura.
Expo añade la cadena de entrega: compilación en la nube, distribución de versiones de prueba y actualizaciones remotas sin pasar por las tiendas \parencite{expo-docs}.
La única fricción conocida (el empaquetador de React Native frente a los enlaces simbólicos de pnpm) tiene su contingencia preparada en la configuración del \eng{monorepo} (sección \ref{sec:alt-monorepo}).

\section{Servicios externos: enfoque decidido, proveedor abierto}
\label{sec:alt-servicios}

Varias funciones dependen de servicios de terceros.
En todas ellas el análisis distingue dos niveles: el \emph{enfoque}, decidido y condicionante del diseño, y el \emph{proveedor} concreto, que se elige en el momento de la integración con dos criterios fijos: capa gratuita suficiente allí donde exista (riesgo R9) y, cuando trate datos personales, garantías europeas de protección de datos (riesgo R7).
Cuando un servicio puede no ofrecer capa gratuita utilizable —el caso del vídeo y de la \ac{IA}—, esa dependencia de pago es justamente lo que relega su función a objetivo secundario (capítulo \ref{ch:objetivos}), con versiones simuladas como contingencia (riesgo R9).

\paragraph{Estimación de plazos y afinidad con \ac{IA}.}
El enfoque decidido es dirigido por eventos: el modelo de lenguaje (\ac{LLM}) se llama únicamente cuando la persona crea o edita una lista o pide una estimación; el sondeo periódico se descarta por costoso e innecesario.
La llamada se ejecuta como trabajo en segundo plano en una cola sobre Redis (BullMQ \parencite{bullmq-docs}) para no bloquear la petición, con respuesta estructurada en formato \ac{JSON} y con caché: si la lista no cambió, no se recalcula.
Las representaciones vectoriales de metas e intereses se generan al crearlos y se guardan en pgvector (sección \ref{sec:alt-bd}).
El proveedor del modelo queda abierto entre la \ac{API} de Anthropic y la de OpenAI; se elegirá al integrar la función comparando el coste por uso y las garantías europeas de protección de datos, los dos criterios fijados al inicio de esta sección.

\paragraph{Verificación de identidad.}
Construir un reconocimiento propio de documentos queda descartado de plano: se integra un proveedor especializado (candidatos: Stripe Identity, Veriff, Onfido) del que se guarda únicamente el estado de la verificación y una clave de referencia, nunca la imagen del documento, por minimización de datos (\ac{RGPD}).
La integración entra en el alcance en modo de pruebas; la contratación del servicio en producción queda fuera (capítulo \ref{ch:objetivos}).

\paragraph{Mensajería en tiempo real.}
El chat exige un canal bidireccional persistente entre cliente y servidor.
Las candidatas son Socket.IO \parencite{socketio-docs}, gestionado sobre el propio servidor, o un servicio equivalente ya gestionado; se decidirá al abordar la mensajería, con la capa gratuita como criterio.

\paragraph{Imágenes y vídeo.}
Las imágenes irán a un almacenamiento de objetos (Cloudflare R2 o S3, abierto).
El vídeo —objetivo secundario desde el 16/07/2026— irá a un servicio gestionado de vídeo (candidatos: Mux o Cloudflare Stream, abierto).

\paragraph{Notificaciones y anuncios.}
Los recordatorios de metas llegarán por notificaciones \eng{push} (Expo Push o Firebase Cloud Messaging, abierto).
La publicidad no se programa desde cero: el producto prevé una red de anuncios existente (AdMob en móvil; AdSense o Ad Manager en web, abierto) más una tabla propia de publicaciones patrocinadas que la ordenación del \eng{feed} inyecta.
Dentro de este \ac{TFG} solo se implementa la tabla propia de patrocinados: la integración de la red comercial queda fuera del alcance (capítulo \ref{ch:objetivos}).

\section{Entrega: de la integración continua heredada al despliegue continuo}
\label{sec:alt-entrega}

De TrackBites se hereda la cadena de comprobaciones automáticas (\eng{pipeline}) de integración continua en GitHub Actions, con estrategia de fallo rápido —análisis estático, comprobación de tipos y pruebas, con las pruebas de extremo a extremo de Playwright como última barrera—, cuya materialización concreta se describe en el capítulo \ref{ch:solucion}.
El salto nuevo es el despliegue continuo: en TrackBites se omitió deliberadamente, y aquí se completa el ciclo, de modo que todo cambio que pasa el \eng{pipeline} llega solo a producción.
El \eng{backend} se empaqueta con Docker, que sustituye el despliegue manual anterior por entornos reproducibles \parencite{docker-docs}.
Los destinos concretos se eligen al abordar el despliegue: la \ac{API}, entre Render, Railway y Fly.io; la web, entre Vercel y Cloudflare Pages; la aplicación móvil, por la cadena de Expo.
La base de datos es la excepción: su destino ya está decidido en Neon (sección \ref{sec:alt-bd}).
El criterio para cerrar los abiertos vuelve a ser el mismo: capa gratuita suficiente para un uso académico.

\section{Resumen del análisis}
\label{sec:alt-resumen}

La tabla \ref{tab:resumen-alternativas} condensa el capítulo: qué se elige, frente a qué alternativa y en qué estado queda cada decisión.

\begingroup
\footnotesize
\setlength{\tabcolsep}{4pt}
\begin{longtable}{@{}>{\raggedright\arraybackslash}p{31mm}>{\raggedright\arraybackslash}p{44mm}>{\raggedright\arraybackslash}p{44mm}>{\raggedright\arraybackslash}p{28mm}@{}}
  \caption{Resumen del análisis de alternativas.}
  \label{tab:resumen-alternativas}\\
  \toprule
  \textbf{Ámbito} & \textbf{Elección} & \textbf{Frente a} & \textbf{Estado} \\
  \midrule
  \endfirsthead
  \multicolumn{4}{c}{\emph{\tablename~\thetable{} — continuación de la página anterior}}\\
  \toprule
  \textbf{Ámbito} & \textbf{Elección} & \textbf{Frente a} & \textbf{Estado} \\
  \midrule
  \endhead
  \midrule
  \multicolumn{4}{r@{}}{\emph{Continúa en la página siguiente…}}\\
  \endfoot
  \bottomrule
  \endlastfoot
  Lenguaje & TypeScript transversal & JavaScript (TrackBites) & Decidida \\
  \addlinespace
  Repositorio & \eng{Monorepo} con \eng{workspaces} de pnpm & Turborepo; repositorios separados* & Decidida \\
  \addlinespace
  Base de datos & PostgreSQL + PostGIS + pgvector & MongoDB + Mongoose; motores especializados aparte & Decidida \\
  \addlinespace
  Base de datos en producción & Neon (región europea) & Supabase; Railway & Decidida \\
  \addlinespace
  Acceso a datos & Prisma (geo/vector en \acs{SQL} crudo) & Mongoose & Decidida \\
  \addlinespace
  Servidor de la \acs{API} & Express 5 estructurado a mano & NestJS & Decidida \\
  \addlinespace
  Pruebas de la \acs{API} & Vitest + Supertest & Jest + Supertest (TrackBites) & Decidida \\
  \addlinespace
  Registro de sucesos & pino + pino-http & winston & Decidida \\
  \addlinespace
  Plataforma \eng{backend} & \eng{Backend} propio (Node + Prisma + PostgreSQL) & Supabase como plataforma integral & Decidida \\
  \addlinespace
  Validación y tipos & Zod compartido en \texttt{packages/core} & Joi solo en el servidor & Decidida \\
  \addlinespace
  Cliente web & React + Vite + TanStack Query & Caché manual con Redux (TrackBites) & Decidida \\
  \addlinespace
  Estado global de interfaz & Pendiente (Redux Toolkit + Reselect en consideración) & — & Abierta \\
  \addlinespace
  Cliente móvil & React Native + Expo & Desarrollo nativo por plataforma* & Decidida \\
  \addlinespace
  \acs{IA} de plazos y afinidad & Por eventos + cola en segundo plano & Sondeo periódico & Enfoque decidido; proveedor abierto \\
  \addlinespace
  Verificación de identidad & Proveedor externo, en modo de pruebas & Reconocimiento propio & Enfoque decidido; proveedor abierto \\
  \addlinespace
  Mensajería en tiempo real & Canal bidireccional persistente & — & Enfoque decidido; proveedor abierto \\
  \addlinespace
  Imágenes y vídeo & Almacén de objetos + servicio de vídeo gestionado & — & Enfoque decidido; proveedor abierto \\
  \addlinespace
  Notificaciones \eng{push} & Expo Push o Firebase Cloud Messaging & — & Enfoque decidido; proveedor abierto \\
  \addlinespace
  Publicidad & Tabla propia de patrocinados (red existente prevista) & Publicidad programada desde cero & Enfoque decidido; integración real fuera del alcance \\
  \addlinespace
  Entrega & Integración continua de fallo rápido + despliegue continuo + Docker & Omitir el despliegue continuo (TrackBites) & Decidida; destinos abiertos \\
  \addlinespace
  \multicolumn{4}{@{}l@{}}{\emph{* Alternativa analizada en este capítulo; no consta como barajada en la documentación previa del proyecto.}} \\
\end{longtable}
\endgroup
% si procede (procede)
% =========================================================================
% Capítulo 7 — Análisis de riesgos (si procede — para BUKIT procede)
% Identificación de riesgos (técnicos, de planificación, legales como
% GDPR/verificación de identidad...), su probabilidad e impacto, y el
% plan de mitigación/contingencia de cada uno.
% =========================================================================

\chapter{Análisis de riesgos}
\label{ch:riesgos}

\todonote{Contenido pendiente.}
     % si procede (procede)
% =========================================================================
% Capítulo 8 — Selección/Descripción de la solución propuesta (si procede)
% Conclusión del análisis de alternativas: la solución elegida, su
% descripción general y por qué. El diseño va en el capítulo 9; el diseño
% de detalle/bajo nivel, al Anexo II.
% =========================================================================

\chapter{Descripción de la solución propuesta}
\label{ch:solucion}

\todonote{Contenido pendiente.}
   % si procede (procede)
% =========================================================================
% Capítulo 9 — Diseño (básico o de alto nivel) (si procede)
% Arquitectura general del monorepo (diagrama), apps (web, móvil, API),
% componentes y sus responsabilidades, modelo de datos, servicios
% externos... El diseño de detalle/bajo nivel va al Anexo II.
% El modelo de datos crece entidad a entidad (diseño incremental, cap. 8):
% cada entidad se documenta aquí al acordarse y construirse.
% =========================================================================

\chapter{Diseño de alto nivel}
\label{ch:diseno}

Este capítulo recoge el diseño del sistema conforme a su método de construcción (capítulo \ref{ch:solucion}): el modelo de datos se presenta entidad a entidad, en el orden en que cada una se acuerda y se implementa, con su porqué y las alternativas descartadas.
\todonote{Pendientes: figura de la arquitectura, componentes de las aplicaciones y contratos de la API, conforme avance el desarrollo; y las figuras UML del modelo de datos y de las secuencias de sesión, ya dibujadas, cuando se decida con el tutor si van en este capítulo o en el Anexo II.}

\section{Modelo de datos}
\label{sec:diseno-modelo}

El modelo vive en un único esquema declarativo del que se derivan tanto las migraciones versionadas de la base de datos como el cliente de acceso a datos tipado que utiliza la \ac{API}: una sola fuente de verdad, el mismo principio que rige la validación compartida (capítulo \ref{ch:alternativas}).
Antes de escribirse, cada entidad pasa una revisión de minimización de datos (riesgo R7, capítulo \ref{ch:riesgos}): cada dato personal debe tener una finalidad declarada, y lo que hoy no la tiene no se almacena.

\subsection{Usuario}
\label{sec:diseno-usuario}

La primera entidad es la cuenta de usuario: identificador, email, hash de la contraseña, nombre de usuario, nivel y fechas de creación y actualización.

El identificador es un \ac{UUID}, por dos razones: no es enumerable (nadie puede recorrer los perfiles incrementando un número ni deducir el volumen de cuentas del sistema) y es un tipo nativo de PostgreSQL.
Se descartaron el entero autoincremental —más compacto y rápido en las uniones, pero enumerable— y el identificador textual por defecto de la capa de acceso a datos (cuid), que se almacena como texto sin tipo nativo.
El criterio queda fijado para el resto de entidades del modelo.
Con la misma vocación de criterio general se fijan, tras la revisión de la entidad: los nombres de la base de datos en minúsculas con guion bajo (hay consultas \acs{SQL} crudas comprometidas para el mapa y la afinidad, y los identificadores con mayúsculas exigirían comillas exactas en cada una), las fechas siempre con zona horaria (los plazos de las metas y las caducidades del producto no pueden depender de la zona horaria de la sesión del servidor), las longitudes máximas respaldadas también en la columna como defensa en profundidad, y los identificadores generados igualmente en la base de datos, no solo en el cliente.

El email y el nombre de usuario son únicos, y ambos se normalizan a minúsculas en el esquema de validación compartido antes de llegar a la base de datos: dos grafías del mismo email o del mismo nombre son la misma identidad.
Esa normalización vive una única vez en la capa de validación compartida, por la que pasa toda escritura (capítulo \ref{ch:alternativas}); se descartó duplicarla en la base de datos con un índice funcional, que la herramienta de migraciones no sabe expresar de forma declarativa y quedaría desincronizado del esquema.
El nombre de usuario existe desde el alta, y no en una edición posterior del perfil, porque el perfil es central en el producto (capítulo \ref{ch:contexto}): cada cuenta nace con su identidad pública.
De la contraseña se guarda únicamente el hash bcrypt, calculado con la biblioteca bcryptjs (capítulo \ref{ch:alternativas}); la contraseña en claro no se persiste ni se registra jamás, y la forma pública de la respuesta se valida contra el contrato compartido, de modo que el hash no puede salir por accidente.
El nivel de usuario (normal o TOP) es un enumerado con normal por defecto y sin ningún punto de la \ac{API} que lo modifique: la admisión TOP llega por un proceso externo (capítulo \ref{ch:contexto}).

La revisión de minimización de la entidad concluye sin hallazgos: no se almacenan teléfono, fecha de nacimiento ni ubicación, y los datos de la verificación de identidad llegarán con su módulo, reducidos al estado y a la clave de referencia del proveedor (capítulo \ref{ch:alternativas}).

\subsection{Sesiones: los \eng{tokens} de refresco}
\label{sec:diseno-sesiones}

La segunda entidad guarda los \eng{tokens} de refresco que emite el inicio de sesión.
Hasta entonces, el \eng{token} de refresco era autosuficiente: valía durante sus treinta días sin que el sistema pudiera anularlo, de modo que no existía un cierre de sesión real ni forma de reaccionar ante un robo.
Guardar cada \eng{token} emitido permite tres cosas: rotarlo en cada uso, cerrar una sesión de verdad y detectar cuándo alguien presenta una copia.

Cada \eng{token} se guarda con su identificador, su usuario, la sesión a la que pertenece, su huella, su fecha de caducidad, la de creación, la de anulación si la hay y el \eng{token} que lo sustituyó al renovarse.
Del \eng{token} solo se guarda su huella SHA-256, nunca el propio \eng{token}: un volcado de la base de datos no entregaría ninguna sesión utilizable.
La huella es SHA-256 y no bcrypt, a diferencia de las contraseñas, porque el \eng{token} ya es una cadena larga y aleatoria: bcrypt es lento a propósito para dificultar que se adivinen contraseñas débiles, y aquí no hay nada débil que proteger, mientras que la renovación sí debe ser rápida.
La caducidad guardada es exactamente la del propio \eng{token}, de modo que una fila caducada corresponde siempre a un \eng{token} que la verificación ya rechaza por sí sola.

La renovación sigue la rotación con detección de reutilización que recomienda la práctica de seguridad vigente para este tipo de \eng{tokens} \parencite[apartado~4.14]{rfc-9700}: cada renovación anula el \eng{token} presentado y entrega uno nuevo, enlazado con el anterior.
Si llega un \eng{token} ya renovado, alguien conserva una copia: el cliente legítimo o un atacante, sin que el sistema pueda saber cuál.
La respuesta es anular la sesión a la que pertenece, con el coste asumido de que el cliente legítimo tenga que volver a iniciar sesión.
Para acotar ese alcance, cada \eng{token} lleva un identificador de sesión que comparten un inicio de sesión y todas sus renovaciones: ante una reutilización se anula solo esa sesión, y las del mismo usuario en otros dispositivos siguen vivas.
Se descartaron anular todas las sesiones del usuario, más simple pero que expulsa también a sus dispositivos legítimos, y una entidad de sesión independiente, más normalizada pero que añadía una tabla sin otra función por ahora.
Dos renovaciones simultáneas con el mismo \eng{token} son indistinguibles de un robo y se tratan igual: la anulación del \eng{token} es condicional dentro de una transacción, de modo que solo una de las dos puede completarse.
Todos los rechazos devuelven la misma respuesta, sin distinguir si el \eng{token} caducó, se cerró o se reutilizó, porque esa distinción solo orientaría a quien prueba \eng{tokens} robados.
Cerrar sesión anula la sesión entera del \eng{token} presentado y responde igual aunque ya estuviera cerrada.

Los \eng{tokens} caducados se borran al iniciar sesión, solo los del propio usuario: no hace falta ninguna pieza nueva y el coste es una consulta por índice.
Se descartaron la purga de toda la tabla en cada inicio de sesión, que añadía una escritura global a la ruta más usada; una tarea programada, que exigía otro proceso con acceso a la base de producción; y no borrarlos, que acumulaba datos personales sin finalidad.

Esta entidad introduce la primera relación del modelo, y con ella el criterio de borrado: los datos que dependen de otra entidad y no tienen valor propio se borran en cascada con ella, decidiéndolo relación a relación.
Las sesiones cumplen ese criterio, así que borrar una cuenta, en ejercicio del derecho de supresión del \ac{RGPD}, se lleva sus sesiones sin ninguna limpieza manual.
Se descartó impedir el borrado mientras existan sesiones y limpiarlas aparte, que complicaba la baja de una cuenta sin aportar nada.
Las relaciones con valor social, como las publicaciones o los seguimientos, decidirán su propio criterio cuando lleguen.

La revisión de minimización de la entidad concluye sin hallazgos: no se guardan la dirección IP, el dispositivo ni el navegador desde los que se inicia sesión.
Serían útiles para una pantalla de sesiones activas, pero ninguna función comprometida los necesita, y son datos personales.
               % si procede (procede)
% =========================================================================
% Capítulo 10 — Descripción de tareas, fases, equipos o procedimientos y
% diagrama de Gantt/cronograma (si procede)
% Metodología de trabajo (decisión SCRUM vs. tradicional/RUP, justificada)
% y las tareas/fases del desarrollo con su cronograma.
% OJO: el Gantt del TFG en sí (obligatorio) tiene su propio capítulo (13).
% =========================================================================

\chapter{Descripción de tareas, fases y procedimientos y diagrama de Gantt}
\label{ch:tareas}

Este capítulo justifica la metodología con la que se planifica y ejecuta el trabajo, define cómo se estima el esfuerzo y presenta el \eng{Product Backlog} inicial y las fases del desarrollo con su cronograma.
La representación gráfica del cronograma completo del \ac{TFG} (incluida la memoria y los hitos académicos) se recoge en el capítulo \ref{ch:gantt-tfg}.

\section{Metodología de trabajo}
\label{sec:metodologia}

\subsection{Restricciones que la metodología debe satisfacer}
\label{sec:metodologia-restricciones}

La elección no parte de preferencias abstractas, sino de las condiciones reales de este proyecto:

\begin{itemize}
  \item \textbf{Alcance emergente.}
        El modelo de datos y las entidades se diseñan de forma incremental y meditada durante el desarrollo, una decisión ya tomada del proyecto; una red completa de tareas definida por adelantado quedaría obsoleta en semanas.
  \item \textbf{Disponibilidad variable y conocida.}
        La dedicación de septiembre será sensiblemente menor por motivos personales (riesgo R2, capítulo \ref{ch:riesgos}); la metodología debe absorber esa variación sin romperse.
  \item \textbf{Una única persona.}
        No hay roles que repartir ni reuniones que sincronizar; toda ceremonia debe justificar su coste ante un equipo de uno.
  \item \textbf{Dos entregables simultáneos.}
        Producto y memoria evolucionan a la vez, así que la redacción debe ser trabajo planificado de primera clase, no un resto.
  \item \textbf{Telemetría para esta memoria.}
        El método debe producir por sí mismo los datos de seguimiento (horas, velocidad, desviaciones) que alimentan este capítulo y el \ref{ch:gantt-tfg}.
  \item \textbf{Marco académico.}
        Unas 360 horas de dedicación de referencia (12 créditos \ac{ECTS}) en un periodo de ejecución objetivo de unos tres meses (capítulo \ref{ch:objetivos}).
\end{itemize}

\subsection{Alternativas consideradas y descartadas}
\label{sec:metodologia-alternativas}

\paragraph{Metodologías predictivas (cascada, \acs{RUP}) con planificación por precedencias.}
El enfoque tradicional visto en la asignatura de Gestión de Proyectos (ciclos tipo \ac{RUP} y técnicas de planificación temporal como \ac{PERT}, Roy o los diagramas de precedencias) exige descomponer el proyecto completo en tareas con duraciones y dependencias estables antes de empezar \parencite{kruchten-rup}.
Esa premisa falla aquí dos veces: el alcance es emergente por decisión de diseño, y la capacidad disponible varía a lo largo del periodo, de modo que la red de tareas exigiría re-planificación continua, un coste permanente sin valor añadido.
A ello se suma que \ac{RUP} prescribe una disciplina de artefactos y roles pensada para equipos, desproporcionada para una persona.
De este enfoque se rescata, no obstante, una pieza: un plan macro de fases e hitos con fechas (capítulo \ref{ch:gantt-tfg}), que da la visión de conjunto que una pila de historias no ofrece.

\paragraph{Scrum puro.}
Scrum fue la metodología empleada por el autor en el proyecto grupal de la asignatura, con resultados satisfactorios, y por ello fue la primera candidata.
Sin embargo, Scrum se define explícitamente sobre un equipo con roles diferenciados (\eng{Product Owner}, \eng{Scrum Master}, \eng{Developers}) y ceremonias cuyo propósito es sincronizar y proteger a ese equipo \parencite{scrum-guide}.
En solitario, la mitad del marco queda vacío: no hay a quién hacer de contrapeso en la negociación del alcance, la \eng{daily} no tiene interlocutor y la \eng{review} no tiene público.
Además, el compromiso cerrado de cada \eng{sprint} encaja mal con una capacidad que se sabe variable: un mes de baja dedicación encadenaría «fracasos» de \eng{sprint} que no informan de nada.
Declarar «se usó Scrum» sería, en rigor, falso; lo honesto es declarar qué elementos de Scrum se adoptan y por qué.

\paragraph{Kanban puro.}
El método Kanban (flujo continuo, trabajo visualizado, límite de trabajo en curso) encaja de forma natural con la disponibilidad variable: si una semana rinde menos, fluye menos trabajo, sin ningún compromiso roto \parencite{anderson-kanban}.
Se descarta como método único porque carece de cadencia: sin iteraciones no hay ritmo fijo de re-planificación ni retrospectiva, y no se genera la medida periódica de velocidad que este \ac{TFG} necesita como telemetría para la memoria y para corregir su primera experiencia de estimación en solitario (riesgo R3).

\subsection{Decisión: método híbrido (Scrumban) adaptado a una persona}
\label{sec:metodologia-decision}

Se adopta un híbrido de la familia Scrumban \parencite{ladas-scrumban, kniberg-scrumban}: la estructura y la telemetría de Scrum que sobreviven al contexto unipersonal, sobre una mecánica de flujo Kanban que absorbe la variabilidad.
La tabla \ref{tab:metodologia-elementos} resume qué se toma de cada marco y con qué adaptación.

\begin{table}[htbp]
  \centering
  \small
  \begin{tabular}{p{3.4cm}p{1.6cm}p{8.2cm}}
    \toprule
    \textbf{Elemento} & \textbf{Origen} & \textbf{Adaptación a una persona} \\
    \midrule
    \eng{Product Backlog} & Scrum & Historias de usuario priorizadas con MoSCoW; toda idea nueva entra por aquí (control del alcance, R1). \\
    Iteraciones & Scrum & \eng{Sprints} semanales (lunes--domingo) con capacidad variable declarada por sprint. \\
    Naturaleza del compromiso & Kanban & El compromiso es el orden del \eng{backlog} y el límite de \ac{WIP}, no un contrato cerrado de \eng{sprint}; el contenido semanal es una proyección. \\
    Límite de \acs{WIP} & Kanban & Máximo 2 elementos en curso: una historia de producto y, como mucho, la épica transversal de memoria; nunca dos historias de producto a medias. \\
    Planificación de \eng{sprint} & Scrum & Sesión individual de unos 20 minutos cada lunes: fijar capacidad y seleccionar del \eng{backlog}. \\
    \eng{Daily} & Scrum & Sustituida por el procedimiento ya vigente de abrir cada sesión releyendo la última entrada del diario de desarrollo. \\
    \eng{Review} & Scrum & Verificación del incremento contra la Definición de Hecho (sección \ref{sec:dod}). \\
    Retrospectiva & Scrum & Sección fija de la entrada del diario al cierre del \eng{sprint}, que incluye la revisión de riesgos (capítulo \ref{ch:riesgos}). \\
    Velocidad & Scrum & Puntos completados por \eng{sprint}; recalibra la estimación de forma continua. \\
    \bottomrule
  \end{tabular}
  \caption{Elementos del método híbrido y su adaptación al contexto unipersonal.}
  \label{tab:metodologia-elementos}
\end{table}

La duración de una semana por \eng{sprint} (frente a las dos semanas empleadas en el proyecto grupal) se justifica por tres razones.
Primera: la realimentación de velocidad llega al doble de frecuencia justo cuando más falta hace, en las primeras semanas de la primera estimación en solitario del autor.
Segunda: el grano semanal permite modular la capacidad de septiembre sin tocar el método (cambia la capacidad, no el ritmo).
Tercera: el argumento clásico contra los \eng{sprints} cortos, el coste proporcional de las ceremonias, desaparece en solitario, donde todas las ceremonias juntas cuestan menos de una hora semanal.

\subsection{Herramienta de gestión}
\label{sec:metodologia-herramienta}

La gestión del \eng{backlog} y del tablero se realiza con GitHub Projects, con las historias como \eng{issues}, los puntos y la épica como etiquetas y las fases como \eng{milestones}.
El tablero tiene cuatro columnas (\eng{Backlog}, Preparado, En curso con límite de \acs{WIP} 2, y Hecho).
Se descarta OpenProject, la herramienta empleada en el proyecto grupal, porque su valor diferencial son las funciones de equipo (roles, actas, flujos de aprobación) que aquí no se usarían, y exige desplegar y mantener un servidor propio.
GitHub Projects, en cambio, vive junto al repositorio y enlaza cada historia con sus \eng{commits} y \eng{pull requests}, lo que genera trazabilidad historia--código sin esfuerzo adicional: evidencia gratuita para esta memoria.
El tablero es la herramienta operativa del día a día; esta memoria conserva las instantáneas relevantes (el \eng{backlog} inicial aquí, el estado final y las desviaciones en los capítulos \ref{ch:resultados} y \ref{ch:gantt-tfg}).

\section{Estimación del esfuerzo}
\label{sec:estimacion}

El esfuerzo se estima en puntos de historia (\eng{story points}), una medida relativa de tamaño, con la escala de Fibonacci modificada (1, 2, 3, 5, 8 y 13).
La estimación relativa se elige frente a la estimación directa en horas porque comparar («esto es como aquello, pero el doble») es más fiable que predecir duraciones absolutas, y porque desacopla el tamaño del trabajo de la capacidad disponible, que aquí varía por diseño \parencite{cohn-estimating}.
Al no existir equipo, no procede el \eng{planning poker}; en su lugar, cada historia se estima por triangulación contra historias de referencia ya conocidas.
La referencia inicial procede del precedente técnico del autor (el proyecto TrackBites): 2 puntos equivalen a «un \eng{endpoint} CRUD sencillo con validación compartida y pruebas de integración», y 1 punto a una tarea mecánica bien conocida.
Toda historia estimada en 13 puntos se considera demasiado grande y debe dividirse antes de entrar en un \eng{sprint} (mitigación del riesgo R3).

Como hipótesis inicial de velocidad se toman unos 11 puntos por \eng{sprint} a plena capacidad (30 horas semanales), es decir, unas 2,7 horas por punto.
Esta equivalencia sirve únicamente para el encaje macro del plan; a partir del tercer \eng{sprint} la manda la velocidad real medida, y la hipótesis se recalibra en cada retrospectiva.

\section{Definición de Hecho}
\label{sec:dod}

Una historia solo se considera terminada cuando cumple todas las condiciones siguientes:

\begin{itemize}
  \item El análisis estático y la comprobación de tipos del \eng{monorepo} pasan en verde.
  \item Las pruebas de la historia existen y pasan, y la suite completa no sufre regresiones.
  \item El cambio está integrado en la rama principal.
  \item Las decisiones tomadas quedan anotadas en el diario de desarrollo, con la alternativa descartada y su porqué.
  \item Si la historia afecta a un capítulo de esta memoria, el capítulo queda actualizado (mitigación del riesgo R4).
\end{itemize}

Esta lista es el criterio de la \eng{review} de cada incremento y protege los objetivos de calidad y trazabilidad del capítulo \ref{ch:objetivos}.

\section{\eng{Product Backlog} inicial}
\label{sec:backlog}

El \eng{backlog} inicial se construye descomponiendo las reglas de producto de \bukit en épicas, y estas en historias de usuario con el formato «como \emph{rol}, quiero \emph{acción}, para \emph{beneficio}».
Cada épica recibe una prioridad MoSCoW \parencite{clegg-moscow}, que traduce el alcance del capítulo \ref{ch:objetivos} a decisiones operativas: \emph{Must} para el alcance comprometido del que dependen los criterios de éxito, \emph{Should} para lo incluido en el alcance sin criterio de éxito asociado, \emph{Could} para los objetivos declarados secundarios, y \emph{Won't} para lo explícitamente excluido (sección \ref{sec:alcance-fuera}), que no se repite aquí.

\begin{table}[htbp]
  \centering
  \small
  \begin{tabular}{clccc}
    \toprule
    \textbf{Id} & \textbf{Épica} & \textbf{Prioridad} & \textbf{Puntos} & \textbf{Fase} \\
    \midrule
    EP01 & Fundaciones técnicas e integración continua      & Must   & 8  & F1 \\
    EP02 & Cuentas y autenticación                          & Must   & 13 & F1 \\
    EP03 & Andamiaje del cliente web y sesión               & Must   & 8  & F1 \\
    EP04 & Listas y metas                                   & Must   & 13 & F2 \\
    EP05 & Publicaciones ligadas a listas                   & Must   & 13 & F2 \\
    EP06 & Grafo social y \eng{feed} de seguidos            & Must   & 10 & F2 \\
    EP07 & Perfil y verificación de identidad               & Must   & 8  & F3 \\
    EP08 & \eng{Feed} personalizado y buscador              & Must   & 13 & F4 \\
    EP09 & Mapa social                                      & Must   & 10 & F3 \\
    EP10 & Afinidad entre personas usuarias                 & Should & 8  & --- \\
    EP11 & Mensajería y notificaciones                      & Must   & 13 & F3 \\
    EP12 & Pruebas \acs{E2E} y despliegue continuo          & Must   & 8  & F4 \\
    EP13 & \acs{IA} de estimación de plazos                 & Could  & 8  & --- \\
    EP14 & Aplicación móvil parcial                         & Could  & 13 & --- \\
    EP15 & Memoria y gestión del proyecto (transversal)     & Must   & 21 & Todas \\
    \bottomrule
  \end{tabular}
  \caption{Épicas del \eng{Product Backlog} inicial, con estimación gruesa y fase prevista.}
  \label{tab:backlog-epicas}
\end{table}

Las estimaciones de la tabla \ref{tab:backlog-epicas} son deliberadamente gruesas: siguiendo la elaboración progresiva propia de los métodos ágiles, solo se detallan en historias las épicas de la fase inminente, y el resto se refina justo antes de abordarse, con el conocimiento acumulado hasta entonces.
Como muestra del nivel de detalle al que se baja cada historia antes de entrar en un \eng{sprint}, se desarrolla una completa:

\begin{quote}
  \textbf{US-11 (EP04, 5 puntos).}
  Como persona usuaria registrada, quiero crear una lista con metas y fecha límite, para convertir mis deseos en compromisos medibles.
  \emph{Criterios de aceptación:} (1) una cuenta gratuita no puede superar las tres listas y el exceso se rechaza con un error explícito; (2) la validación vive una sola vez en el paquete compartido y la usan tanto la \ac{API} como el cliente web; (3) toda lista exige al menos una meta y una fecha límite; (4) el flujo queda cubierto por pruebas de integración.
\end{quote}

La tabla \ref{tab:backlog-f1} recoge el desglose completo de la primera fase, ya refinado.

\begin{table}[htbp]
  \centering
  \small
  \begin{tabular}{clcc}
    \toprule
    \textbf{Id} & \textbf{Historia (resumen)} & \textbf{Épica} & \textbf{Puntos} \\
    \midrule
    US-01 & Cadena de integración continua: lint, tipos y pruebas     & EP01 & 3 \\
    US-02 & Manejo centralizado de errores y registro de sucesos      & EP01 & 3 \\
    US-03 & Despliegue del esqueleto de la \acs{API}                  & EP01 & 2 \\
    US-04 & Modelo de datos de usuario y registro de cuentas          & EP02 & 3 \\
    US-05 & Inicio de sesión con \acs{JWT} de acceso y refresco       & EP02 & 3 \\
    US-06 & Rotación y revocación de \eng{refresh tokens}             & EP02 & 5 \\
    US-07 & Doble capa de limitación de peticiones                    & EP02 & 2 \\
    US-08 & Andamiaje web con enrutado de las cuatro pestañas         & EP03 & 3 \\
    US-09 & Cliente \acs{API} tipado y gestión de sesión en la web    & EP03 & 3 \\
    US-10 & Pantallas de registro e inicio de sesión                  & EP03 & 2 \\
    \bottomrule
  \end{tabular}
  \caption{Desglose en historias de la fase F1 (fundaciones), 29 puntos.}
  \label{tab:backlog-f1}
\end{table}

\section{Fases y cronograma del desarrollo}
\label{sec:fases}

El periodo de ejecución se organiza en cuatro fases que agrupan trece \eng{sprints} semanales, del 20/07/2026 al 18/10/2026, precedidas por una fase de preparación ya en curso.
Cada fase cierra con un incremento demostrable, que actúa como hito de control.

\begin{table}[htbp]
  \centering
  \small
  \begin{tabular}{clllc}
    \toprule
    \textbf{Fase} & \textbf{Contenido} & \textbf{Sprints} & \textbf{Fechas} & \textbf{Horas} \\
    \midrule
    F0 & Preparación: entorno, memoria base, este plan & ---        & hasta 19/07      & $\sim$30 \\
    F1 & Fundaciones: EP01, EP02, EP03                 & S1--S3     & 20/07 -- 09/08   & 90 \\
    F2 & Núcleo social: EP04, EP05, EP06               & S4--S6     & 10/08 -- 30/08   & 90 \\
    F3 & Diferenciadores: EP07, EP09, EP11             & S7--S11    & 31/08 -- 04/10   & 70 \\
    F4 & Cierre de producto: EP08, EP12                & S12--S13   & 05/10 -- 18/10   & 60 \\
    \bottomrule
  \end{tabular}
  \caption{Fases del desarrollo con su capacidad estimada (la épica transversal EP15 se reparte entre todas).}
  \label{tab:fases}
\end{table}

La capacidad de cada \eng{sprint} refleja la disponibilidad real esperada: 30 horas semanales en verano y octubre, y 10 horas semanales en septiembre (\eng{sprints} S7 a S10), conforme al riesgo R2 ya identificado.
Por eso la fase F3 abarca cinco \eng{sprints}: su duración en semanas no indica más trabajo, sino menos capacidad por semana.
\todonote{Confirmar las capacidades estimadas por sprint (30/10/30 horas) y la fecha de inicio de S1 antes del 20/07.}

\begin{table}[htbp]
  \centering
  \small
  \begin{tabular}{cllc}
    \toprule
    \textbf{Sprint} & \textbf{Fechas} & \textbf{Tema orientativo} & \textbf{Horas} \\
    \midrule
    S1  & 20/07 -- 26/07 & Integración continua mínima; modelo de usuario y registro & 30 \\
    S2  & 27/07 -- 02/08 & Autenticación completa con rotación de \eng{tokens}       & 30 \\
    S3  & 03/08 -- 09/08 & Andamiaje web y flujo de sesión; \emph{cierre de F1}      & 30 \\
    S4  & 10/08 -- 16/08 & Listas y metas                                            & 30 \\
    S5  & 17/08 -- 23/08 & Publicaciones desde listas                                & 30 \\
    S6  & 24/08 -- 30/08 & Grafo social y \eng{feed} de seguidos; \emph{cierre de F2} & 30 \\
    S7  & 31/08 -- 06/09 & Perfil y verificación de identidad                        & 10 \\
    S8  & 07/09 -- 13/09 & Perfil y verificación de identidad (continuación)         & 10 \\
    S9  & 14/09 -- 20/09 & Mapa social                                               & 10 \\
    S10 & 21/09 -- 27/09 & Mapa social (continuación)                                & 10 \\
    S11 & 28/09 -- 04/10 & Mensajería y notificaciones; \emph{cierre de F3}          & 30 \\
    S12 & 05/10 -- 11/10 & \eng{Feed} personalizado y buscador                       & 30 \\
    S13 & 12/10 -- 18/10 & \acs{E2E} y despliegue continuo; \emph{cierre de F4}      & 30 \\
    \bottomrule
  \end{tabular}
  \caption{Calendario de \eng{sprints} con su capacidad; el tema es una proyección, el compromiso real es el orden del \eng{backlog}.}
  \label{tab:sprints}
\end{table}

\subsection{Encaje del plan y tensión declarada}
\label{sec:fases-tension}

El inventario comprometido (\emph{Must}) suma 138 puntos, mientras que la capacidad nominal de los trece \eng{sprints} (310 horas, a la hipótesis de 2,7 horas por punto) da unos 115 puntos.
Existe por tanto una tensión inicial de unos 23 puntos ($\approx$ 60 horas, unas dos semanas a plena capacidad), que se declara aquí en lugar de maquillarse ajustando estimaciones a la baja.
La decisión consciente es triple: se mantiene el 18/10/2026 como fecha objetivo ambiciosa, se comunica una horquilla de cierre realista entre esa fecha y mediados de noviembre de 2026 (capítulo \ref{ch:gantt-tfg}), y las prioridades \emph{Should} y \emph{Could} solo entran si la velocidad real supera la hipótesis.
Esta tensión es, de hecho, la primera observación registrada del riesgo R3, y su evolución se vigila desde el primer cierre de \eng{sprint}.
En cualquier caso, la holgura estructural del calendario académico (defensa en junio de \thesisYear) hace que esta tensión amenace al objetivo propio del autor, no al \ac{TFG}.
 % si procede (procede)
% =========================================================================
% Capítulo 11 — Cálculos, algoritmos (si procede)
% Para BUKIT: algoritmo de ranking del feed (ponderación TOP/normal),
% affinity matching con pgvector, estimación de plazos con IA...
% Si finalmente no procede, quitar su \input de main.tex y borrar este
% archivo.
% =========================================================================

\chapter{Cálculos y algoritmos}
\label{ch:algoritmos}

\todonote{Contenido pendiente.}
  % si procede
% =========================================================================
% Capítulo 12 — Descripción/Análisis de resultados (si procede)
% Qué se ha conseguido: funcionalidades completadas, métricas (cobertura
% de tests, rendimiento...), capturas del producto final.
% Si finalmente no procede, quitar su \input de main.tex y borrar este
% archivo.
% =========================================================================

\chapter{Análisis de resultados}
\label{ch:resultados}

\todonote{Contenido pendiente.}
  % si procede
% =========================================================================
% Capítulo 13 — Diagrama de Gantt del TFG (OBLIGATORIO)
% Planificación del propio TFG (no del producto): fases del trabajo
% académico, cronograma previsto vs. real. Fuente de los datos reales:
% el dev-journal (docs/dev-journal/, con las horas de cada sesión).
% =========================================================================

\chapter{Diagrama de Gantt del TFG}
\label{ch:gantt-tfg}

\todonote{Contenido pendiente.}
            % obligatorio
% =========================================================================
% Capítulo 14 — Descripción del presupuesto/gastos del TFG (OBLIGATORIO)
% Horas de trabajo (valoradas), licencias, infraestructura (hosting,
% dominios, APIs de IA...). La guía permite añadir apartados: un análisis
% de rentabilidad (modelo de negocio: publicidad, cuentas TOP...) sumaría.
% =========================================================================

\chapter{Descripción del presupuesto y gastos del TFG}
\label{ch:presupuesto}

\todonote{Contenido pendiente.}

% Apartado opcional (añadido permitido por la guía) — descomentar si aplica:
% \section{Análisis de rentabilidad}
% \label{sec:rentabilidad}
          % obligatorio
% =========================================================================
% Capítulo 15 — Conclusiones (OBLIGATORIO)
% Resultados obtenidos a lo largo del trabajo, valoración de las
% conclusiones alcanzadas, y cuestiones pendientes que darían lugar a una
% profundización posterior (prospectiva / trabajo futuro).
% =========================================================================

\chapter{Conclusiones}
\label{ch:conclusiones}

\todonote{Contenido pendiente. Se escribe al final, junto con el resumen. Incluir la lección del 12/11/2026 sobre la asistencia de IA: complicó el seguimiento del proyecto más de lo debido (exceso de reglas, pérdida de contexto y de coherencia), salvo en las guías breves, y obligó a simplificar el método (preliminar «Empleo de IA», capítulos 7 y 10).}
         % obligatorio

% ----- BIBLIOGRAFÍA (normas APA) -----
\printbibliography[heading=bibintoc]

% ----- ANEXOS -----
\appendix
% La plantilla numera los anexos en romanos: Anexo I, II, III (no A, B, C).
\renewcommand{\thechapter}{\Roman{chapter}}
% =========================================================================
% Anexo I — Pliego de condiciones. Normativa aplicable. Plan de pruebas
% Normativa: GDPR/LOPDGDD (datos personales, geolocalización, KYC),
% requisitos de edad (18+ para descubrimiento), etc.
% Plan de pruebas: estrategia de testing (unit/integración/E2E) y CI.
% =========================================================================

\chapter{Pliego de condiciones. Normativa aplicable. Plan de pruebas}
\label{ap:pliego}

\todonote{Contenido pendiente.}

% =========================================================================
% Anexo II — Diseño de detalle
% Esquemas, diagramas de bajo nivel, modelo de datos completo (Prisma),
% fragmentos de código relevantes... según la plantilla del centro.
% =========================================================================

\chapter{Diseño de detalle}
\label{ap:diseno-detalle}

\todonote{Contenido pendiente.}

% =========================================================================
% Anexo III — Manuales de usuario y de administrador
% Si procede: guía de uso de la app y guía de despliegue/administración.
% =========================================================================

\chapter{Manuales de usuario y de administrador}
\label{ap:manuales}

\todonote{Contenido pendiente.}


\end{document}
